\section{Future Work}
\label{sec:future}

\paragraph{L3 evaluation.}
The collapse mechanism (Section~\ref{sec:l3}) is implemented but not
yet evaluated at scale. Open questions include: how much context does
model-authored compaction recover compared to system-initiated
eviction? Do declared losses in collapse summaries enable accurate
fault decisions (can the model determine from a summary alone whether
it needs the original content)? What is the quality impact of
replacing 20~turns of dialogue scaffolding with a single outcome
sentence? These require longitudinal measurement across sessions of
varying length and task complexity.

\paragraph{Cost-aware eviction pressure.}
The current eviction policy uses a space-based threshold (context
window capacity). The inverted cost model argues for cost-based
pressure: 23\% memory pressure sounds low, but 45{,}000 tokens per
turn at inference pricing is real money. Eviction pressure should be
denominated in token-turns (the product of page size and projected
remaining lifetime), not in fractional capacity. This reframes
eviction from ``are we running out of space?'' to ``is this page
earning its keep?''

\paragraph{Cost-weighted pin decay.}
Pinned pages should not remain resident forever. A fault tells us
the content was needed \emph{then}, not forever. We propose
cost-weighted decay: pin strength halves every $K$~turns since last
access. A page becomes evictable when its projected keep cost exceeds
its fault cost. Under low memory pressure, this is permissive
(pins persist). Under high pressure, cold pins are released first.
This gives LRU-like behavior with cost weighting, and prevents the
monotonic working-set growth that permanent pinning causes in long
sessions.

\paragraph{Phase-aware eviction.}
Planning and execution have different working set characteristics.
Planning requires broad simultaneous context (many files held at
once); execution is sequential (read, edit, move on). The proxy
could detect phase from request patterns---many Reads with no Edits
suggests planning---and adjust the eviction threshold accordingly.

\paragraph{Per-connection isolation.}
The current implementation uses a single PageStore for all
connections. Subagent sessions share eviction state with the main
session, causing cross-contamination. Per-connection isolation would
give each session its own eviction history, fault tracking, and
pinning state.

\paragraph{Trace-driven simulation.}
We extract reference strings (the sequence of page accesses) from
session transcripts and replay them under different replacement
policies. This enables evaluation of policies against recorded
workloads without live inference-priced experiments. The inverted cost
model (Section~\ref{sec:inverted-cost}) means the optimal offline
policy is not Belady's MIN but a cost-weighted variant that minimizes
total keep cost plus fault cost. We plan to derive and evaluate this
cost-optimal offline bound alongside MIN for comparison.

\paragraph{Cross-session access pattern prediction.}
The break-even formula requires estimating $T_{\text{until\_next\_ref}}$---the
number of turns until a page is next referenced. Within a single
session, this can only be estimated from local patterns. Across
sessions, the prediction becomes tractable: coding sessions exhibit
regular access patterns (orientation reads of project structure,
iterative edits of a working file, periodic reference to configuration).
A Markov model trained on cross-session reference strings could
predict re-reference probability per page per turn, enabling
cost-optimal eviction decisions. The data collection is nearly free:
the proxy already logs every eviction, every fault, every tool call.
What is needed is a persistent, queryable store of access patterns
that survives session boundaries---precisely the role of a cross-session
memory system.

\paragraph{Cooperative extensions.}
Phantom tools open a side-channel between the model and the proxy.
The current vocabulary---release and fault---is minimal. Extensions
include: \emph{prefetch hints} (model signals upcoming working set
needs), \emph{eviction classes} (groups of pages evicted as a unit),
\emph{variable-fidelity tombstones} (AST skeletons instead of bare
path markers), and \emph{ephemeral flags} (tool results marked for
immediate eviction after extraction). Each expands the model's
ability to communicate cognitive state to the memory manager.

\paragraph{Beyond paging: object-addressed memory.}
The current system manages fixed-granularity pages addressed by file
path. But the natural unit for LLM memory management is the semantic
object---a conversation phase, a design decision, a debugging session.
These objects vary in size, have relationships to each other, and can
be compressed to multiple fidelity levels rather than the binary
resident/evicted state of hardware pages. The eviction tombstone
becomes a compressed summary that functions as both a page table
entry (retrieval handle to the backing store) and a cache line
(enough semantic content to answer queries without faulting).
Declared losses in the summary tell the model what it \emph{cannot}
answer from the summary alone, guiding the decision to fault.
The backing store is queryable---the model can ask a question of
evicted content without materializing it in full. This shifts the
abstraction from block-addressed paging to object-addressed memory
management, where eviction is cooperative, compression is authored,
and the backing store is a database rather than a swap partition.

\paragraph{Comparative measurement.}
Instrument other agentic tools (Cursor, Copilot, Windsurf,
Continue.dev) to measure the same metrics. The structural argument
predicts similar waste profiles; empirical confirmation would
strengthen generalizability.

% ====================================================================
